\chapter{Conclusion}

The present work investigates whether the same concept of copying the solution of an ordinary differential equations system, as was done for the Lorenz-1960 set, can be applied to the simple feedforward artificial neural network with a much lower accuracy of the reference data, which is in the form of numerical data. Motivation is practical and academic tension. There are classical numerical solvers like Runge-Kutta methods that provide dependable reference solutions, but they solve the system step by step on a fixed grid. A continuous input/output mapping can be learned over time by neural approaches. Although PINNs have been introduced for the solution of differential equations in recent efforts, state variables are still relevant classes of variables. Losses, or operator-learning architectures \cite{matthews2024pinnde, babni2025error, lau2020oden}. This study thus focuses on narrows. The data-driven feedforward ANN is trained to approximate $(x(t), y(t), z(t))$ for the Lorenz-1960 initial value problem, without adding ODE residual terms to the loss function.

The research design is controlled architecture comparison with validated numerical baseline. The parameters for the Lorenz-1960 formulation, initial conditions and time. The system used in the reference problem of Chapter 2 (PinnDE) is used in the intervals. Two numerical solvers are used to produce the reference trajectory and to validate it. The reference trajectory is generated and validated with a custom numerical solver, as well as a commercial numerical solver. RK4 is applied before using it as training data and SciPy DOP853 is used. The ANN study is divided into two parts. In phase 1 the depth, width and optimizer are kept constant and the activation function is changed. During the first phase the depth, width and optimizer are fixed and the activation function is changed. Phase 2 compares Adam, L-BFGS and SGD with momentum on the top-performing phase 1 architectures. This separation is made to make the interpretation of the architecture and optimiser effects easier. The predicted result of the work will be an empirical benchmark to make an approximation of the Lorenz-1960 system to a simple ANN. Previous works have shown that neural networks are able to solve, or to approximate, differential equations of a wide range of formulations \cite{mall2013comparison, okereke2021novel, dubey2024solving}. The physics-informed approaches have been used for Lorenz-type problems  \cite{matthews2024pinnde, aslam2024neuro}. The academic merit of this project is that it is based on the simpler and more straightforward question of how far a standard feedforward ANN can be pushed by systematic variations in its architecture and an exploration of the results it learns. A verified numerical solution is called a training target. The solution can help determine if the bracelet is or is not better. What about the physics informed constraints, is it required for this benchmark or can plain-ANN design be used? These are still necessary to get a higher accuracy and reliability. At this report stage the supported contribution is: The validated baseline design, the dataset and training protocol, and the planned evaluation framework. If the experiment records the declaration of a final best ANN architecture then the report does not state a final best ANN architecture. In Chapter 4, you will find metric tables and plots. This is intentional. Without supporting evidence, any claim of accuracy would detract from the study. Final conclusions regarding the optimal depth, width, activation function or optimizer should be finalized only after careful experimental research and analysis. Planned comparison is completed and research needs to be continued in the study. Three directions. First, the most suitable configuration needs to be tested repeatedly over random Use of seeds to check for stability. Second, in order to assess the generalization of the trained model, it needs to be tested on denser Use time grids or other initial conditions. Third, the ANN results obtained in the plain should be compared with the physics-based, literature informed and operator learning baselines in a more direct manner. These extensions would assist in deciding if the architecture chosen for this project is effective for only one or more of the objects, controlled trajectory or for wider neural approximation of nonlinear ODE systems.
